% !TeX root = ../papers/nc-p3.tex
\section{特異基点の吹上げと全ての整な階数一五次曲線の排除}
\label{sec:rank-one-singular-quintics}
\begin{theorem}\label{thm:rank-one-singular-quintics}
四辺有限の標準有限Frobenius窓で入力は被約、反対平方根は一般可逆で全二次式を成分に持つとする。
階数一軌跡のどの既約曲線成分も次数4以下である。
\end{theorem}
\begin{proof}[二次曲面と算術種数]
既に総次数は5以下である。整五次曲線 $C$ があれば非平面的で算術種数は2以下である。
滑らかな場合と、正規化が有理で二次曲面に載らない場合は前定理による。
算術種数が1以下で特異なら正規化は有理である。滑らかな二次曲面上の整五次の二次数は
$(1,4)$ または $(2,3)$ で算術種数は0か2である。二次錐でも整五次の算術種数は2である。
実際解消 $\mathbb F_2$ の例外曲線を $E$、超平面を $H=E+2f$ とすれば真変換は $bE+5f$ で $1\le b\le2$。
錐の次数 $t$ の切断で曲線上消えるものは解消上の $tH-bE-5f$ の切断である。
$E$ との負交叉により固定 $E$ を除くと、$t\ge3$ で可動部分は $(t-3)H+f$、切断数は $(t-2)(t-1)$。
従ってHilbert多項式は $(t+1)^2-(t-2)(t-1)=5t-1$ である。

残る算術種数2の特異曲線でも、Riemann--Rochにより $h^0(L)=h^0(L')=4$, $h^0(2L)=h^0(LL')=9$ である。
双対化層は階数一捩れなし、負次数に非零切断はなく非特殊性が従う。
両埋込みは完全で唯一の二次曲面 $Q,Q'$ に載る。定理\ref{thm:rank-one-genus-two-quintic}の
二次式比較で $L'=L\eta$, $\eta$ は非自明な二捩れ、曲線外の印数 $s\ge7$ である。
曲線外の各印は両側の対応する二次曲面の外にある。同定理の一点を通る一般平面族は曲線の有限な特異点を避け、
五つの共役点と一点での平方の関係の証明がそのまま成り立つ。
\end{proof}
\begin{proof}[割線の長さと滑らかな固定点]
曲線外の二印の割線の切断長を出力で $b$ とする。$b\le2$、$b=0$ は先と同じ一般五点で除かれる。
$b=2$ の方向は一つの曲線外点 $p'$ から高々四つで、いずれも他の印を高々一つ持つ。
実際 $p'$ からの平面射影 $C\to\Gamma$ は平面五次への有限双有理射で、
$\mathcal O_\Gamma\subset f_*\mathcal O_C$ の余核長は算術種数の差 $6-2=4$ である。
長さ二の繊維は少なくとも1を使い、三つの曲線外印と長さ二因子を含む直線は、前の四次小行列式により
階数一軌跡の直線成分になってしまう。

また入力割線が $C$ を長さ二で切れば同じ二印の出力割線もその長さ二部分を含む。
入力直線上のその部分と二印の和は長さ4で、二次評価を消す完全跡を持つ。
これは直線上の長さ4因子の随伴、またはアフィン一変数環で最高次の三次係数を取れば分かる。
二乗座標の同型で出力一次座標の制限はこの四次元完全対合（単元倍を許す）で全等方、従って次元2以下である。
よって出力でもスキームとして二印の割線上にある。

$b=1$ の固定点 $r$ が滑らかなら、先と同じ一般四点の鉛筆の移送により入力射影次数は4である。
入力固定部分は滑らかな一点 $z$ である。正規化で固定次数は1で、特異点を基点にすれば
一般平面の局所交数は重複度2以上になるからである。
$L(-z)=L'(-r)$ は曲線自体の束の等式である。両鉛筆は滑らかな基点を取った後は全点で射となり、
一般点での射影同値が全体に延びるからである。従って $r,z$ は異なる滑らかな二点で $2r\sim2z$。
次数二の基点なし鉛筆があれば一意である。異なる二つがあれば $C$ が $(2,2)$ 曲線に有限双有理となり、
像の算術種数1が $C$ の算術種数2より小さくなってしまう。ゆえに $r,z$ はその一意な鉛筆の全分岐点であり
（各点の分岐指数が2）、固定した $\eta$ が許す順序なし二点組は高々一つである。
二組の比較 $z+r'\sim z'+r$ は前定理と同一で、異なる有効因子ならその次数二鉛筆の繊維となる。
これらの点を通る $b=1$ の割線に三つの外部印がないことも先と同じである。
\end{proof}
\begin{proof}[特異固定点の局所長二を残す]
$b=1$ の固定点 $r$ が特異とする。上の $C\to\Gamma$ のここでの繊維長は1で、有限環拡大は
中山の補題により $1$ で生成される。よって射はここで局所同型、$r$ は平面曲線特異点である。
三重以上なら局所デルタが少なくとも3となり $p_a(C)=2$ を超えるので二重点である。
割線を通る一般平面が $r$ に切る部分 $T$ は長さ2で、割線との交わりは長さ1である。
残りの三点は正規化からの次数三鉛筆の共役点である。$T$ と二つの固定印の平方は独立である。
実際割線上の異なる三点は二次で独立であり、割線の式と $r$ で消えない一次式との積で
$T$ の付加的な接方向を検出できる。

一般平面内の長さ5の切断に対する二次式は依然 $Q'$ の制限が唯一である（線型正規性の切断完全列）。
従って二つの外部印を加えた平方の関係は両印で非零、先の独立性から三つの共役点でも全て非零である。
$T$ における関係の汎関数は、その積対合の零化イデアルを除くと長さ $j=0,1,2$ の完全対合となる。
入力一次座標の全等方性はこの商を取った後にも成り立つ。
$j=0$ では長さ5の完全対合上で評価は二次元以下となり、入力二印を結ぶ固定直線に動点が入ってしまう。
$j=2$ では七次元で評価は三次元以下、入力平面が出力の全長さ五因子を含み $L=L'$ となる。
よって $j=1$ であり、二印・$r$ の対応点・三動点が入力で共面である。
平面が一定なら $C$ が退化するので、入力割線は $r$ の対応点を通る。
入力側の切断長も1である。長さ2なら直前の四次元跡の議論で出力切断が長さ2になってしまう。

従って両鉛筆の基点イデアルは抽象曲線上で全く同じ $\mathfrak m_r$ である。
$g:C_1\to C$ をそこでの吹上げとする。両鉛筆は $C_1$ 上で $g^*L\otimes\mathfrak m_r\mathcal O_{C_1}$ と
$g^*L'\otimes\mathfrak m_r\mathcal O_{C_1}$ の基点なし二切断となる。
三動点上での共面性により入力射は出力の次数三射を経由し、入力次数も5以下の正の三の倍数なので3である。
従って二射は底の射影同値を除き一致し、上の束の比較から $g^*\eta=1$ を得る。
$C_1$ は算術種数1の整曲線である。平面二重点の極大イデアルでの吹上げは局所平面の吹上げでの真変換であり、
随伴で算術種数の減少は $2(2-1)/2=1$ となる（他の点は同型）。
\end{proof}
\begin{proof}[二つの種類の固定点は共存しない]
$\eta$ が一点の吹上げで自明になる特異点は高々一つである。異なる二点の各々で自明化できるなら、正規化上の
非零定数倍まで一意の自明化がそれぞれの点の外で $C$ に降りる。両開集合を合わせれば $\eta=1$ となる。
またそのような特異固定点があれば滑らかな $b=1$ の固定点は存在しない。後者があれば
$\eta=\mathcal O(r'-z)$ であり、二つの異なる滑らかな点の差が $C_1$ で自明となる。
その単純極と単純零を持つ関数は算術種数1の整曲線を射影直線に次数1で有限に写してしまう。

特異固定点がない場合、各印が結べる他の外部印は $b=2$ に四つ、滑らかな $b=1$ に二つまでである。
$s\ge7$ との等号では、指定した一方の滑らかな固定点を通る割線上で七印が二個ずつ組になる矛盾を得る。
特異固定点が一つなら、どの外部印からも $b=1$ で使える直線はその点へ向かう一本だけで、
直線切断長上界3により他の外部印はその上に二つ以下である。再び全てが等号で七印となるが、
今度はその固定点を通る各直線がちょうど三個ずつ外部印を分ける。これも不可能である。
\end{proof}
